% Unit 134 terjemahan Bahasa Indonesia dari chapter9.tex baris 1950--2035.
% Sumber: Methods of Algebra, Volume 2, commit
% 9a5803ff77dd3257484cb177f851a73770a59dd3 (CC BY 4.0).
% stable-unit-id: o014.aljabr2.chapter9.exercises
% Cakupan: seluruh Latihan Bab 9.

% segment-id: o014.li2.chapter9.exercises.h001
\begin{Exercises}
	% segment-id: o014.li2.chapter9.exercises.ex001
	\item Misalkan $\Bbbk$ medan dan $V,W$ ruang vektor-$\Bbbk$; seperti
	biasa, definisikan $V^\vee:=\Hom_{\Bbbk}(V,\Bbbk)$, dan seterusnya.
	Tuliskan pemetaan linear kanonik
	$V^\vee\dotimes{\Bbbk}W^\vee\to
	(V\dotimes{\Bbbk}W)^\vee$. Tunjukkan bahwa pemetaan ini selalu injektif,
	tetapi tidak surjektif bila $V$ dan $W$ keduanya berdimensi tak hingga.

	% segment-id: o014.li2.chapter9.exercises.ex002
	\item Misalkan $T$ sembarang objek dari kategori monoidal $\mathcal{C}$,
	dan $\iota:U\to\munit$ suatu morfisme. Dengan asumsi ${}^*U$ (atau
	$U^*$) ada, buktikan bahwa bijeksi dalam
	Proposisi~\ref{prop:dual-internal-Hom}(iii)
	% segment-id: o014.li2.chapter9.exercises.q001
	\[ \Hom(T\otimes U,T)\to\Hom(T,T\otimes{}^*U) \]
	(atau $\Hom(U\otimes T,T)\to\Hom(T,U^*\otimes T)$) memetakan
	$\identity_T\otimes\iota$ ke $\identity_T\otimes{}^*\iota$ (atau
	memetakan $\iota\otimes\identity_T$ ke
	$\iota^*\otimes\identity_T$).

	% segment-id: o014.li2.chapter9.exercises.ex003
	\item Misalkan kategori monoidal $\mathcal{C}$ juga kategori aditif,
	dengan objek nol $0$, dan andaikan $\otimes$ bifunktor aditif. Tunjukkan
	bahwa $X\otimes0=0=0\otimes X$ untuk setiap $X$.

	% segment-id: o014.li2.chapter9.exercises.ex004
	\item Misalkan $\mathcal{C}$ kategori monoidal beranyam yang kaku,
	sekaligus kategori abelian, dan $\End_{\mathcal{C}}(\munit)$ suatu medan.
	Untuk $X\in\Obj(\mathcal{C})$ beserta dual kanannya $X^*$, buktikan
	bahwa $X\neq0$ ekuivalen dengan $X\otimes X^*\neq0$, dan juga ekuivalen
	dengan $X^*\otimes X\neq0$.
	% segment-id: o014.li2.chapter9.exercises.hint001
	\begin{hint}
		Perhatikan bahwa $X\neq0$ ekuivalen dengan
		$\mathrm{ev}_X\neq0$, dan juga dengan $\mathrm{coev}_X\neq0$;
		terapkan Proposisi~\ref{prop:rigid-unit-simple}.
	\end{hint}

	% segment-id: o014.li2.chapter9.exercises.ex005
	\item Misalkan $\mathcal{C}$ dan $\mathcal{C}'$ kategori monoidal
	simetris yang kaku, keduanya juga kategori abelian, dengan
	$\End_{\mathcal{C}}(\munit)$ medan dan $\munit'\neq0$. Buktikan bahwa
	setiap funktor monoidal eksak yang kompatibel dengan struktur kepang,
	$F:\mathcal{C}\to\mathcal{C}'$, memenuhi
	$X\neq0\iff F(X)\neq0$.
	% segment-id: o014.li2.chapter9.exercises.hint002
	\begin{hint}
		Ulangi argumen Lema~\ref{prop:Tannaka-generalities}(i), atau gunakan
		metode latihan sebelumnya.
	\end{hint}

	% segment-id: o014.li2.chapter9.exercises.ex006
	\item (Nobuo Yoneda) Misalkan $\mathcal{C},\mathcal{D}$ kategori dan
	$F:\mathcal{C}^{\opp}\times\mathcal{C}\to\mathcal{D}$ funktor. Sebuah
	\emph{baji} untuk $F$ berarti data
	$\left(W,(e_X)_{X\in\Obj(\mathcal{C})}\right)$, yang juga ditulis
	$e:W\xrightarrow{\bullet}F$, dengan
	\begin{compactitem}
		\item $W$ objek $\mathcal{D}$;
		\item $e_X:W\to F(X,X)$ morfisme dalam $\mathcal{D}$,
	\end{compactitem}
	sedemikian rupa sehingga diagram berikut komutatif untuk setiap
	$f:X\to Y$ dalam $\mathcal{C}$:
	% segment-id: o014.li2.chapter9.exercises.q002
	\[\begin{tikzcd}
		W \arrow[r, "e_Y"] \arrow[d, "e_X"'] & F(Y, Y) \arrow[d, "{F(f, Y)}"] \\
		F(X, X) \arrow[r, "{F(X, f)}"' inner sep=0.6em] & F(X, Y).
	\end{tikzcd}\]

	Secara dual, definisikan \emph{kobaji} sebagai data $(M,(f_X)_X)$,
	yang ditulis $f:F\xrightarrow{\bullet}M$, dengan syarat komutativitas
	% segment-id: o014.li2.chapter9.exercises.q003
	\[\begin{tikzcd}
		F(Y, X) \arrow[r, "{F(f, X)}"] \arrow[d, "{F(Y, f)}"'] & F(X, X) \arrow[d, "{f_X}"] \\
		F(Y, Y) \arrow[r, "{f_Y}"'] & M.
	\end{tikzcd}\]
	\begin{enumerate}[(i)]
		\item Jelaskan bagaimana semua baji (atau kobaji) untuk $F$ membentuk
		kategori. Jika kategori baji (atau kobaji) memiliki objek terminal
		(atau awal), objek itu disebut \emph{end} (atau \emph{koend}) dari
		$F$, dengan notasi
		% segment-id: o014.li2.chapter9.exercises.q004
		\[ \text{end}\;=\int_{X\in\mathcal{C}}F(X,X),\qquad
		\text{koend}\;=\int^{X\in\mathcal{C}}F(X,X). \]
		Integral di sini hanyalah notasi, tetapi bukan tanpa alasan.
		\index{end@end (end)}
		\index{koend@koend (coend)}
		\index[sym1]{integral@$\int F(X, X)$}
		\item Tafsirkan \eqref{eqn:cogebra-L} dan konstruksi
		$\coHom(\omega_1,\omega_2)$ dalam
		Proposisi~\ref{prop:cogebra-L-universal} sebagai kasus khusus koend.
		\item Andaikan $\mathcal{C}$ kecil secara esensial relatif terhadap
		semesta Grothendieck yang dipilih
		(Definisi~\ref{def:essentially-small-cat}), dan $\mathcal{D}$ lengkap
		(atau kokomplet). Buktikan bahwa setiap $F$ memiliki end (atau koend).
		\item Untuk dua funktor
		$F,G:\mathcal{C}^{\opp}\times\mathcal{C}
ightrightarrows
		\mathcal{D}$, definisikan \emph{transformasi dinatural}
		$\alpha:F\xrightarrow{\bullet}G$ sebagai data
		$(\alpha_X)_{X\in\Obj(\mathcal{C})}$, dengan
		$\alpha_X:F(X,X)\to G(X,X)$, sedemikian rupa sehingga diagram berikut
		komutatif untuk setiap $f:X\to Y$ dalam $\mathcal{C}$:
		% segment-id: o014.li2.chapter9.exercises.q005
		\[\begin{tikzcd}[row sep=tiny]
			& F(X, X) \arrow[r, "{\alpha_X}"] & G(X, X) \arrow[rd, "{G(X, f)}"] & \\
			F(Y, X) \arrow[ru, "{F(f, X)}"] \arrow[rd, "{F(Y, f)}"'] & & & G(X, Y). \\
			& F(Y, Y) \arrow[r, "{\alpha_Y}"'] & G(Y, Y) \arrow[ru, "{G(f, Y)}"'] &
		\end{tikzcd}\]
		Tunjukkan bahwa baji dan kobaji di atas merupakan kasus khusus
		transformasi dinatural. Tunjukkan pula bahwa transformasi alami, yakni
		morfisme antarfunktor $\mathcal{C}\to\mathcal{D}$ atau
		$\mathcal{C}^{\opp}\to\mathcal{D}$, juga merupakan kasus khususnya.
		\index{transformasi dinatural@transformasi dinatural (dinatural transformation)}
	\end{enumerate}
	Pemaparan rinci mengenai end dan koend tersedia dalam monograf
	\cite{Lor21}.

	% segment-id: o014.li2.chapter9.exercises.ex007
	\item Tinjau funktor $F,G:\mathcal{C}\to\mathcal{D}$, dengan
	$\mathcal{C}$ kecil secara esensial. Buktikan bahwa himpunan morfisme
	antara keduanya dapat ditafsirkan sebagai end:
	% segment-id: o014.li2.chapter9.exercises.q006
	\[ \Hom_{\mathcal{D}^{\mathcal{C}}}(F,G)\simeq
	\int_{X\in\mathcal{C}}\Hom_{\mathcal{D}}(FX,GX). \]
	Dengan ini, tafsirkan pusat $Z(\mathcal{C})$ dari kategori $\mathcal{C}$
	sebagai end yang ditentukan oleh $\Hom_{\mathcal{C}}$.

	% segment-id: o014.li2.chapter9.exercises.ex008
	\item Misalkan $\Bbbk$ medan. Buktikan bahwa koend funktor
	$\Hom_{\Bbbk}:\cate{Vect}_{\mathrm f}(\Bbbk)^{\opp}\times
	\cate{Vect}_{\mathrm f}(\Bbbk)\to\cate{Vect}_{\mathrm f}(\Bbbk)$
	dapat diidentifikasikan dengan $\Bbbk$ beserta pemetaan jejak
	% segment-id: o014.li2.chapter9.exercises.q007
	\[ \Tr_V:\End_{\Bbbk}(V)\to\Bbbk,\qquad
	V:\;\text{ruang vektor-$\Bbbk$ berdimensi hingga}. \]

	% segment-id: o014.li2.chapter9.exercises.ex009
	\item Secara umum, end menyerupai operasi mengambil invarian, sedangkan
	koend menyerupai perekatan. Jika pembaca mengenal himpunan
	simplisial, tafsirkan funktor realisasi geometris $|\mathord\cdot|$ dari
	\S\ref{sec:geom-realization} sebagai koend
	% segment-id: o014.li2.chapter9.exercises.q008
	\[ |X|\simeq\int^{[n]\in\simpDelta}X_n\times|\Delta^n|,qquad
	X\in\Obj(\cate{sSet}). \]

	% segment-id: o014.li2.chapter9.exercises.ex010
	\item Untuk kategori kecil $\mathcal{C}$, definisikan kategori funktor
	$\mathcal{C}^\wedge:=\cate{Set}^{\mathcal{C}^{\opp}}$. Buktikan bahwa
	untuk setiap $\mathcal{F}\in\Obj(\mathcal{C}^\wedge)$ terdapat
	isomorfisme kanonik dalam $\mathcal{C}^\wedge$
	% segment-id: o014.li2.chapter9.exercises.q009
	\[ \int^{X\in\Obj(\mathcal{C})}\mathcal{F}(X)\times
	\Hom_{\mathcal{C}}(\cdot,X)\rightiso\mathcal{F}. \]
	Koend ini berasal dari funktor
	% segment-id: o014.li2.chapter9.exercises.q010
	\[ \mathcal{C}^{\opp}\times\mathcal{C}\to\mathcal{C}^\wedge,
	\quad(X,Y)\mapsto\mathcal{F}(X)\times\Hom_{\mathcal{C}}(\cdot,Y). \]
	% segment-id: o014.li2.chapter9.exercises.hint003
	\begin{hint}
		Lihat \S\ref{sec:Yoneda}, khususnya Teorema Kepadatan
		\ref{prop:Yoneda-density}.
	\end{hint}
	\index{pembenaman Yoneda!kepadatan (density)}

	% segment-id: o014.li2.chapter9.exercises.ex011
	\item Untuk kategori Tannakian $\mathcal{T}$, funktor serat
	$\omega:\mathcal{T}\to\cated{Mod}B$, dan aljabar-$B$ komutatif $B'$,
	tunjukkan bahwa, di bawah isomorfisme
	Korolari~\ref{prop:Aut-Tannakian}, unsur identitas dan operasi invers
	dalam grup
	$\Aut^{\otimes}(\omega\dotimes{B}B')$ masing-masing berasal dari
	kounit dan antipoda $\coEnd_B(\omega)$ melalui
	$\Hom_{B\dcate{CAlg}}(\mathord\cdot,B')$.
	% segment-id: o014.li2.chapter9.exercises.hint004
	\begin{hint}
		Identitas dan invers suatu grup ditentukan secara unik oleh
		perkalian. Perkaliannya telah diketahui bersesuaian dengan
		komultiplikasi $\coEnd(\omega)$.
	\end{hint}

	% segment-id: o014.li2.chapter9.exercises.ex012
	\item Misalkan $G$ grup, $G\dcate{Set}$ kategori semua himpunan kecil
	dengan aksi kiri $G$, dan $U:G\dcate{Set}\to\cate{Set}$ funktor lupa.
	Berikan isomorfisme grup alami
	% segment-id: o014.li2.chapter9.exercises.q011
	\[ G\simeq\Aut(U). \]
	Jelaskan bagaimana ini dapat dipandang sebagai prototipe sederhana
	Teorema Tannaka--Krein~\ref{prop:Tannaka-Krein}.
	% segment-id: o014.li2.chapter9.exercises.hint005
	\begin{hint}
		Terapkan Lema Yoneda pada
		$U\simeq\Hom_{G\dcate{Set}}(G,\mathord\cdot)$ untuk memperoleh
		isomorfisme monoid $G\simeq\End(U)$.
	\end{hint}
\end{Exercises}
